\documentclass[10pt,a4paper]{amsart}
\usepackage[margin=19mm]{geometry}
\usepackage{amsmath,amssymb,mathtools}
\usepackage[scr=rsfs,cal=euler]{mathalfa}
\usepackage{xfrac}
\usepackage[unicode,hidelinks,bookmarks=false]{hyperref}
\newcommand{\ie}{i.e.\ }
\newcommand{\cf}{cf.\ }
\newcommand{\resp}{respectively\ }
\newcommand{\Subsection}{\subsection}
\providecommand{\nobreakdash}{\nobreak\hbox{-}}
\setlength{\emergencystretch}{2em}
\multlinegap0pt
\usepackage{bm}
\usepackage{bbm}
\usepackage{dsfont}
\usepackage{xspace}
\usepackage{cancel}
\usepackage{mathtools}
\usepackage{amsthm}
\makeatletter
\@ifundefined{theorem}{\newtheorem{theorem}{Theorem}}{}
\@ifundefined{proposition}{\newtheorem{proposition}[theorem]{Proposition}}{}
\makeatother
\title[The Isoperimetric Problem in Regular Trees]{The Isoperimetric Problem in Regular Trees}
\author{Marc Troyanov}
\date{Source: arXiv:2604.18769v1 (CC BY 4.0)}
\begin{document}
\maketitle

\noindent\emph{Source and licence.} The Isoperimetric Problem in Regular Trees. Version arXiv:2604.18769v1 (CC BY 4.0), distributed under the Creative Commons Attribution 4.0 International licence, as recorded in the arXiv licence metadata for this exact version and shown on the abstract page https://arxiv.org/abs/2604.18769v1. Full rights notice: licenses/arxiv-2604.18769-CC-BY-4.0.txt.\par\smallskip

\noindent\textit{Editorial scope.} Counted unit: the Proposition whose source label is \texttt{prop:tau-gluing} in the article (source0.tex lines 1089--1098, source label \texttt{prop:tau-gluing}), delivered together with the article's own complete original proof of it (source0.tex lines 1100--1124, one \texttt{proof} environment of 25 lines). No mathematics is written, repaired or re-derived: statement and proof are the article's own text line for line. Required context retained uncounted: eq.degree-sum (lines 1076--1078). Every definition, display and label of those lines is retained unchanged. Omitted from this package and not redistributed: 1--1075, 1079--1088, 1099--1099, 1125--2096. Nothing inside a retained excerpt is omitted or altered: every retained body is delivered exactly as the article's lines are written, so no omission receipt of any kind accompanies this package. Citations: the retained excerpts carry no citation command at all, so no bibliography entry of the article is transcribed and no reference is redistributed. Reference adaptation: none was needed and none was made; every reference inside the delivered unit resolves inside it, so no unresolved reference or citation is produced. Printed slips kept literal: no printed word, symbol, hypothesis or number of the counted statement or of its proof was altered. Layout and class: the article's own class is \texttt{article} with packages including geometry, amsmath, amsthm, amssymb, amsfonts, enumerate, hyperref, mathrsfs, float, caption, subcaption, authblk, inputenc, fontenc; the wrapper is the lane's pinned \texttt{amsart} editorial wrapper at 10pt on A4, which restates the theorem-like environments the retained text needs and adds the article's own macro definitions that the retained text needs (none), with extra wrapper packages (bm, bbm, dsfont, xspace, cancel, mathtools). The delivered numbering is the unit's own. Assets: no retained span contains a figure environment, a graphics-inclusion command, a PDF-page inclusion, a table or a TikZ picture; no image file is copied into this package, the package declares no asset dependency and no image file is redistributed with it. Licence: version arXiv:2604.18769v1 of this article is distributed under the Creative Commons Attribution 4.0 International licence, as recorded in the arXiv licence metadata for this exact version and shown on the abstract page https://arxiv.org/abs/2604.18769v1. A full rights notice is delivered at licenses/arxiv-2604.18769-CC-BY-4.0.txt. Characters: every retained excerpt is pure ASCII, so the delivered engine, XeLaTeX with its default Latin Modern text font, reports no missing character.
\begin{equation}\label{eq.degree-sum}
   \deg_D(p)=\sum_{i=1}^m \deg_{D_i}(p),
\end{equation}

\begin{proposition}\label{prop:tau-gluing}
Let $D$ be obtained by gluing the domains $D_1,\dots,D_m$ at~$p$. Then
\[
   \tau(D) = \begin{cases}
       \displaystyle \sum_{i=1}^m \tau(D_i) + (m-1),          &\text{if } p\in\partial D, \\[0.3cm]
       \displaystyle \sum_{i=1}^m \tau(D_i) - (d-m),           &\text{if } p\notin\partial D .
     \end{cases}
\]
In particular, gluing increases the total branching excess  $\tau$ by exactly  $m-1$ if  $p\in \partial D$, while $\tau$ remains unchanged or decreases if  $p \not\in \partial D$.
\end{proposition}

\begin{proof}
For each $i$ and any vertex  $x\ne p$ lying in $D_i$ we have $\deg_D(x)=\deg_{D_i}(x)$.  As for the point $p$,  from \eqref{eq.degree-sum} we have
\begin{equation}\label{eq.deg-expand}
   \deg_D(p)-1  = \left( \sum_{i=1}^m \deg_{D_i}(p)\right) - 1    = \sum_{i=1}^m (\deg_{D_i}(p)-1) + (m-1).
\end{equation}
We now separately consider the cases where $p$ belongs to the boundary of $D$ or not.

\smallskip 

If  $p\in\partial D$, then  using~\eqref{eq.deg-expand}, we obtain
\begin{align*}
\tau(D)  &= \sum_{x\in \partial D\setminus \{p\}} \bigl(\deg_D(x)-1\bigr)    + \bigl(\deg_D(p)-1\bigr) 
 \\[0.3em] &=  \sum_{i=1}^m    \left(\sum_{x\in \partial D_i\setminus \{p\}}    \bigl(\deg_{D_i}(x)-1\bigr)\right)    + \bigl(\deg_D(p)-1\bigr) \\[0.3em]
&= \sum_{i=1}^m    \left(\sum_{x\in \partial D_i}    \bigl(\deg_{D_i}(x)-1\bigr)\right)    + (m-1)  = \sum_{i=1}^m \tau(D_i) + (m-1).
\end{align*}


\smallskip 

If $p\notin\partial D$, the argument is similar, except that the contribution $\deg_D(p)-1$ must \textit{not} be counted in the computation of $\tau(D)$. Since in this case $\deg_D(p)=d$, we obtain
\[
 \tau(D) = \sum_{i=1}^m \tau(D_i) + (m-1) - (d-1)  = \sum_{i=1}^m \tau(D_i) - (d-m).
\]
This concludes the proof.
\end{proof}

\end{document}
